(i) Work in $\Lambda V$ through $\Phi$ (\claimref{std-algebra.exterior-isomorphism}). If the $v_j$ are dependent, write one as a combination of the others; multilinearity expands $A$ into outer products with a repeated factor, each zero. If they are independent, extend them to a basis $v_1, \dots, v_n$; then $A$ is a member of the basis $\{v_I\}$ of $\Lambda V$, so $A \ne 0$. For $x = \sum_j x_j v_j$, $x \wedge A = \sum_{j > k} x_j\, v_j \wedge A$, a combination of distinct basis elements, so $x \wedge A = 0$ exactly when $x \in \operatorname{span}(v_1, \dots, v_k)$.

(ii) If $[A] = [A']$, write the factors of $A'$ in the factors of $A$ with a $k \times k$ matrix $M$; multilinearity and alternation give $A' = \det(M) A$, and $\det M \ne 0$ as both are nonzero. Proportional blades have the same subspace by (i).

(iii) Let $w_1, \dots, w_m$ be independent with $w_j \wedge Y = 0$, and extend to a basis $w_1, \dots, w_n$ of $V$. Expand $Y = \sum_{|I| = k} y_I w_I$. For $j \le m$, $w_j \wedge w_I$ is $0$ if $j \in I$ and a distinct basis element otherwise, so $y_I = 0$ unless $\{1, \dots, m\} \subseteq I$. With $Y \ne 0$ this forces $m \le k$, and $m = k$ forces $Y = y\, w_1 \wedge \dots \wedge w_k$. Conversely a nonzero blade attains $k$ by (i).

(iv) If the $v_i$ are dependent, both sides vanish: $A = 0$ by (i), and the Gram matrix of dependent vectors is singular. Otherwise $A \ne 0$ and the restriction of $Q$ to $[A]$ has an orthogonal basis $u_1, \dots, u_k$ (\claimref{synthesis.completion-recovers-the-real-clifford-algebras}(iii)); write $v_i = \sum_j M_{ij} u_j$. Then $A = \det(M)\, u_1 \wedge \dots \wedge u_k = \det(M)\, u_1 \cdots u_k$, since orthogonal vectors anticommute and their antisymmetrised product is their product (\claimref{std-algebra.grade-decomposition}). Using $\tilde{A} = \det(M)\, u_k \cdots u_1$ (\claimref{std-algebra.involution-properties}), $A \tilde{A} = \det(M)^{2} \prod_j Q(u_j)$. The Gram matrix of the $v_i$ is $M \operatorname{diag}(Q(u_j)) M^{\mathsf T}$, with the same determinant. That determinant vanishes exactly when the restricted form is degenerate. If it is nonzero, $\tilde{A}/(A\tilde{A})$ is a two-sided inverse, as $\tilde{A} A = (A \tilde{A})^{\sim} = A\tilde{A}$ is a scalar. If it is zero and $A \ne 0$, then $A \tilde{A} = 0$ with $\tilde{A} \ne 0$, so $A$ is a zero divisor and has no inverse.
